\documentclass[11pt]{article}
\usepackage[margin=1in]{geometry}
\usepackage{amsmath,amssymb,amsthm,booktabs}
\usepackage{xcolor}
\usepackage[colorlinks,linkcolor=blue!45!black]{hyperref}

\newcommand{\T}{\mathsf{T}}
\newcommand{\R}{\mathbb{R}}
\newcommand{\spn}{\operatorname{span}}
\newtheorem{proposition}{Proposition}
\newtheorem{lemma}{Lemma}
\newtheorem{corollary}{Corollary}
\newtheorem{measurement}{Measurement}
\theoremstyle{remark}
\newtheorem{remark}{Remark}

\setlength{\parindent}{0pt}
\setlength{\parskip}{5pt}

\begin{document}

\section*{Why the skew condition is the stiff one}

Measurement~4.30 of the derivation document reports that at unit relative violation the
endpoint map responds to a skew violation of hypothesis (i) by $0.647$ and to a symmetric
off-span violation of hypothesis (ii) by $0.0982$, a ratio of $6.585$, and records the value as
measured, not predicted. This note derives the mechanism, computes the constants in closed form
in two tractable regimes, and shows that the ratio is not a universal constant: it diverges as
the marginals homogenize, and its value for the sampled ensemble is set by the ensemble's own
marginal separation.

Throughout, $V(t)=(1-t)^2\Sigma_0+t(1-t)(C+C^{\T})+t^2\Sigma_1$,
$G(t)=(1-t)(C^{\T}-\Sigma_0)+t(\Sigma_1-C)$, $K=GV^{-1}$, $\dot M=KM$, $M(0)=I$, and
$\mathcal D(E)=M_0\int_0^1\Phi(t)^{-1}K_E(t)\Phi(t)\,dt$ is the first-order response of
Proposition~4.28, with
$K_E(t)=[(1-t)E^{\T}-tE]V_0(t)^{-1}-G_0(t)V_0(t)^{-1}[t(1-t)(E+E^{\T})]V_0(t)^{-1}$.
Write $E=S_E+N_E$ for the symmetric and skew parts. The expansion point is any admissible
symmetric coupling, which is hypothesis (i) of Theorem~4.11; hypothesis (ii) is not needed for
the response formula.

\begin{lemma}\label{lem:gv}
For any symmetric coupling $C_0$, $\,G_0(t)=\tfrac12 V_0'(t)$.
\end{lemma}

\begin{proof}
$V_0'=-2(1-t)\Sigma_0+2(1-2t)C_0+2t\Sigma_1
=2\big[(1-t)(C_0-\Sigma_0)+t(\Sigma_1-C_0)\big]=2G_0$.
\end{proof}

\begin{proposition}[homogeneous marginals]\label{prop:homog}
Let $\Sigma_0=\Sigma_1=\Sigma$ and $C_0=\kappa\Sigma$ with $\kappa$ admissible. Then $M_0=I$ and,
for every perturbation $E$,
\[
\mathcal D(S_E)=0,
\qquad
\mathcal D(N_E)=-I_\kappa\,N_E\Sigma^{-1},
\qquad
I_\kappa=\int_0^1\frac{dt}{\omega_\kappa(t)},
\quad
\omega_\kappa(t)=1-2(1-\kappa)t(1-t),
\]
and for $\kappa\in(-1,1)$, with $\mu=2(1-\kappa)\in(0,4)$,
$I_\kappa=\tfrac{4}{\sqrt{\mu(4-\mu)}}\arctan\sqrt{\tfrac{\mu}{4-\mu}}\ \ge1$.
In particular the symmetric response vanishes for \emph{every} symmetric direction, in or out of
$\spn\{\Sigma_0,\Sigma_1\}$, while every skew direction responds at a rate bounded below by
$\|N_E\Sigma^{-1}\|$.
\end{proposition}

\begin{proof}
$V_0=\omega_\kappa\Sigma$ and, by Lemma~\ref{lem:gv}, $G_0=\tfrac12\omega_\kappa'\Sigma$, so
$K_0=\tfrac12(\log\omega_\kappa)'\,I$ is scalar and $\Phi$ is a scalar multiple of the identity;
$\int_0^1(\log\omega_\kappa)'=0$ gives $M_0=I$, and the conjugation by $\Phi$ in $\mathcal D$
cancels, leaving $\mathcal D(E)=\int_0^1K_E\,dt$. For symmetric $E$,
\[
K_E=\Big[\frac{1-2t}{\omega_\kappa}
+(1-\kappa)\frac{2t(1-t)(1-2t)}{\omega_\kappa^{2}}\Big]\,E\Sigma^{-1},
\]
and both coefficients are odd under $t\mapsto1-t$ while $\omega_\kappa$ is even, so the integral
vanishes. For skew $E=N$, $(1-t)E^{\T}-tE=-N$ and $E+E^{\T}=0$, so
$K_E=-N\Sigma^{-1}/\omega_\kappa$ and the second statement follows. The closed form is the
standard integral of the reciprocal quadratic.
\end{proof}

\begin{proposition}[commuting model, pairwise action]\label{prop:pair}
Let $\Sigma_0=\mathrm{diag}(a)$, $\Sigma_1=\mathrm{diag}(b)$, and $C_0=\mathrm{diag}(s)$
admissible, and write $v_i(t)=(1-t)^2a_i+2t(1-t)s_i+t^2b_i$. Then
$\Phi(t)=\mathrm{diag}\sqrt{v_i(t)/a_i}$, $M_0=\mathrm{diag}\sqrt{b_i/a_i}$, and for any
perturbation with pairwise parts $p_{ij}=\tfrac12(E_{ij}+E_{ji})$,
$q_{ij}=\tfrac12(E_{ij}-E_{ji})$,
\[
\mathcal D(E)_{ij}
=\sqrt{\tfrac{b_i}{a_j}}\,\big(p_{ij}\,\mathcal A_{ij}-q_{ij}\,\mathcal B_{ij}\big),
\qquad
\mathcal B_{ij}=\int_0^1\frac{dt}{\sqrt{v_iv_j}}\;>0,
\]
\[
\mathcal A_{ij}
=\int_0^1\frac{(1-2t)-t(1-t)(\log v_i)'}{\sqrt{v_iv_j}}\,dt
\;=\;-\frac12\int_0^1 t(1-t)\,\big(\log\tfrac{v_i}{v_j}\big)'(t)\,
\frac{dt}{\sqrt{v_iv_j}}.
\]
\end{proposition}

\begin{proof}
Lemma~\ref{lem:gv} gives $K_0=\tfrac12 V_0'V_0^{-1}$, diagonal, so
$\varphi_i=\sqrt{v_i/a_i}$ and $M_0$ as stated. Entrywise, with $V_0,G_0$ diagonal,
\[
K_E(t)_{ij}
=\frac{(1-t)E_{ji}-tE_{ij}}{v_j}
-\frac{v_i'}{2v_i}\cdot\frac{2t(1-t)\,p_{ij}}{v_j}
=\Big[(1-2t)-t(1-t)(\log v_i)'\Big]\frac{p_{ij}}{v_j}-\frac{q_{ij}}{v_j},
\]
and $(\Phi^{-1}K_E\Phi)_{ij}=K_{E,ij}\,\varphi_j/\varphi_i$ with
$\varphi_j/\varphi_i=\sqrt{a_iv_j/(a_jv_i)}$; multiplying by $(M_0)_i$ gives the display. For the
second form of $\mathcal A_{ij}$, note
$\big[(1-2t)-t(1-t)(\log v_i)'\big]/v_i=\tfrac{d}{dt}\big[t(1-t)/v_i\big]$, so the first form is
$\int_0^1\sqrt{v_i/v_j}\;\tfrac{d}{dt}\big[t(1-t)/v_i\big]dt$; integrate by parts, the boundary
term vanishes at both ends, and
$\tfrac{d}{dt}\sqrt{v_i/v_j}=\tfrac12\sqrt{v_i/v_j}\,(\log(v_i/v_j))'$.
\end{proof}

\begin{corollary}\label{cor:supp}
$\mathcal A_{ij}=0$ whenever $v_i\equiv\lambda v_j$ for a constant $\lambda$ --- in particular on
the diagonal $i=j$, for homogeneous marginals, and for any pair with
$(b_i,s_i)/a_i=(b_j,s_j)/a_j$. In general
\[
|\mathcal A_{ij}|\;\le\;\frac18\,\sup_{t}\big|\big(\log\tfrac{v_i}{v_j}\big)'(t)\big|\;
\mathcal B_{ij},
\qquad\text{hence}\qquad
\frac{\mathcal B_{ij}}{|\mathcal A_{ij}|}\;\ge\;
\frac{8}{\sup_t\big|\big(\log\tfrac{v_i}{v_j}\big)'\big|}.
\]
\end{corollary}

\begin{remark}\label{rem:mech}
The mechanism is now visible. The skew response integrates the kernel $1/\sqrt{v_iv_j}$, which is
strictly positive on all of $[0,1]$: no cancellation is possible. The symmetric response
integrates the same kernel multiplied by two suppression factors: the window $t(1-t)\le\tfrac14$,
vanishing at both endpoints, and the rate at which the two interpolant variances separate,
$(\log(v_i/v_j))'$, which vanishes identically when the marginals are homogeneous. A symmetric
perturbation moves the endpoint map only through the \emph{variance} channel, which the parity of
the clean flow cancels except insofar as coordinates traverse different variance profiles; a skew
perturbation enters the \emph{drift} directly. The ratio of the two responses is therefore of
order $8/\delta$ with $\delta$ the typical log-variance mismatch rate, which for independent
Wishart marginals is order one --- giving a ratio of several, not of orders of magnitude, and not
a universal constant.
\end{remark}

\begin{measurement}\label{meas:ratio}
All checks in float64 against the same second-order integrator as Measurement~4.30
(\texttt{notebooks/ratio\_proof.py}, register \texttt{runs/ratio/results}).
(1)~The pairwise formulas of Proposition~\ref{prop:pair} match the finite-differenced endpoint
map over $33$ random pair perturbations to maximum relative error $1.3\times10^{-6}$; the two
forms of $\mathcal A_{ij}$ agree to $2.8\times10^{-17}$. (2)~On engineered pairs with
proportional variance profiles the measured symmetric response is below $9.1\times10^{-9}$, as
Corollary~\ref{cor:supp} requires, including off-span directions. (3)~With homogeneous marginals
and $\Sigma$ non-diagonal, so that the diagonal framework of Proposition~\ref{prop:pair} does not
apply, the response to a random symmetric perturbation is $3.5\times10^{-10}$
while the skew response in the same setup is $1.54$, matching
$-I_\kappa N\Sigma^{-1}$ to relative error $3.2\times10^{-9}$ ($\kappa=0.85$,
$I_\kappa=1.0532$). (4)~Interpolating the second marginal toward the first,
$\Sigma_1(\delta)=\Sigma_0+\delta(\Sigma_1-\Sigma_0)$ on the paper's ensemble ($n=16$, $64$
systems, median response per unit violation at relative violation $0.05$):
\begin{center}
\begin{tabular}{@{}lcccccc@{}}
\toprule
$\delta$ & $0.05$ & $0.1$ & $0.2$ & $0.4$ & $0.7$ & $1.0$\\
\midrule
skew/span ratio & $76.8$ & $42.0$ & $24.2$ & $14.6$ & $10.1$ & $7.8$\\
\bottomrule
\end{tabular}
\end{center}
The fitted log--log slope is $-0.756$ with rms residual $0.017$. The first-order prediction for
the $\delta\to0$ limit is slope $-1$: by Corollary~\ref{cor:supp} the span response carries a
factor $\sup_t|(\log v_i/v_j)'|$, which is $O(\delta)$ when the marginals differ by $\delta$,
while Proposition~\ref{prop:homog} makes the skew response $O(1)$ in that limit.

The measured slope over the reported window is milder than $-1$, and it is worth being exact
about why, since the ratio slope is identically the skew slope minus the span slope (verified to
$10^{-15}$). Fitting the two responses separately, and extending the sweep a decade below the
window above:
\begin{center}
\begin{tabular}{@{}lccc@{}}
\toprule
window & span exponent & skew exponent & ratio slope\\
\midrule
$\delta\in[0.05,1]$    & $0.573$ & $-0.159$ & $-0.732$\\
$\delta\in[0.005,0.05]$ & $0.933$ & $-0.019$ & $-0.952$\\
\bottomrule
\end{tabular}
\end{center}
The discrepancy is therefore \emph{not} attributable to the decline of the skew constant: that
decline contributes $-0.16$, which pushes the ratio slope toward $-1$, not away from it. It is
the span response scaling sublinearly, as $\delta^{0.57}$, over a window that is not yet
asymptotic. A decade lower the span exponent has moved to $0.93$ and the ratio slope to $-0.95$,
which is the convergence the first-order argument predicts. Slopes are ensemble-dependent at the
$0.02$ level: the same fit on an independently drawn pair of marginals gives $-0.756$ against
$-0.732$ here. At the ensemble's own
separation ($\delta=1$) the sweep gives $7.8$, against $6.585$ from the regression-intercept
estimator of Measurement~4.30 at dimensions $2$--$64$.
\end{measurement}

\begin{remark}
Consequences. First, the sentence ``measured, not predicted'' should be replaced: the ratio is
the quotient of two response constants whose structure is given by
Propositions~\ref{prop:homog}--\ref{prop:pair}, the span response carrying the double suppression
of Corollary~\ref{cor:supp} and the skew response none. Second, the value $6.585$ is a property
of the sampled ensemble at its own marginal separation; it grows without bound as the marginals
homogenize, so no universal constant should be claimed, and the registered contradiction ``skew
and span equally fragile'' is thereby explained rather than merely recorded. Third, in the
near-homogeneous regime the operative content of Theorem~4.11 is the skew condition alone: the
span condition is soft there, failing only at second order in the marginal mismatch.
\end{remark}

\end{document}
